\input{../tex-inputs/latex-leseansicht-vorspann}

\section[Arthur Schnitzler und Olga Gussmann an Gustav Schwarzkopf, 17. 6. 1901]{L04502 Arthur Schnitzler und Olga Gussmann an Gustav Schwarzkopf, 17. 6. 1901}
\nopagebreak\mylabel{L04502v}
\rehead{ }\normalsize\beginnumbering\briefempfaengerindex{Schwarzkopf, Gustav@\textsc{Schwarzkopf, Gustav}!zzzSchnitzler, Olga@\emph{von Olga Schnitzler}!1901-06-171@{17. 6. 1901}|(be}\briefempfaengerindex{Schwarzkopf, Gustav@\textsc{Schwarzkopf, Gustav}!zzzSchnitzler, Arthur@\emph{von Arthur Schnitzler}!1901-06-171@{17. 6. 1901}|(be}
\toendnotes[C]{\smallbreak\pagebreak[2]}
\correspDesc{Versand  durch Arthur Schnitzler, Olga Gussmann am 17. 6. 1901 in Salzburg
\newline{}Übermittlung  am 17. 6. 1901 in Salzburg
\newline{}Zustellung  am 18. 6. 1901 in I., Innere Stadt
\newline{}Erhalt  durch Gustav Schwarzkopf im Zeitraum [18. 6. 1901 – 22. 6. 1901?] in Wien}\toendnotes[C]{\smallbreak}
\Standort{DLA, A:Schnitzler, HS.1985.1.1897.}
\physDesc{Bildpostkarte, 195 Zeichen
\newline{}Handschrift Arthur Schnitzler: Bleistift, deutsche Kurrent
\newline{}Handschrift Olga Schnitzler: Bleistift, lateinische Kurrent
\newline{}Versand: 1) Stempel: »\nobreak{}\oindex{Salzburg@\textbf{Salzburg}|pwk}Salzburg-Bahnhof, 17. 6. 01, 11 A\nobreak{}«.   2) Stempel: »\nobreak{}\oindex{I., Innere Stadt@\textbf{I., Innere Stadt}|pwk}Wien 1/1 1, 18/6. 1901, 9–10½ V., bestellt\nobreak{}«. }\toendnotes[C]{\smallbreak}
\pstart
           \noindent{}\centering{}{\pb}\textcolor{gray}{\textbf{Residenz-Platz\oindex{Residenzplatz@\textbf{Residenzplatz}|pw} mit Hofbrunnen und
                  Glockenthurm.}}\hspace*{2.5em}\textcolor{gray}{\textbf{Salzburg\oindex{Salzburg@\textbf{Salzburg}|pw}.}}\pend
           \pstart{}{\pb}\textsc{Herrn Gustav
                     Schwarzkopf}\pend{}\pstart{}\textsc{Wien}\oindex{Wien@\textbf{Wien}|pw}\pend{}\pstart{}\textsc{I. Tiefer
                        Graben 23}\oindex{Wien@\textbf{Wien}!I., Innere Stadt@\textbf{I., Innere Stadt}!Tiefer Graben 17@\textbf{Tiefer Graben 17}|pw}\pend{}{\bigskip}\vspace{1em}
\pstart
           \noindent{}Herzliche Grüße. Wir{ }ſind vorläufig hier geblieben wo es{ }ſchön iſt und regnet, was{ }ſich zuweilen miteinander verträgt.\pend
           \pstart Ihr \spacefill\mbox{A. S.}\pend{}\selectlanguage{ngerman}\vspace{1em}
\pstart
           \noindent{}{[}hs. O. Schnitzler:{]} \label{T_L04502-1v}\edtext{Herzlichen Gruss. \spacefill\mbox{O.G.}}{\lemma{\textnormal{\emph{Herzlichen Gruss. O.G.}}}\Cendnote{\textnormal{Olga Gussmanns\pwindex{Schnitzler, Olga 17.\,1.\,1882 Wien – 13.\,1.\,1970 Lugano@\textsc{Schnitzler, Olga} (17.\,1.\,1882 Wien – 13.\,1.\,1970 Lugano)|pwk} Gruß findet sich am
                  oberen Rand der Karte.}}}\label{T_L04502-1}\pend
           \selectlanguage{ngerman}\endnumbering\briefempfaengerindex{Schwarzkopf, Gustav@\textsc{Schwarzkopf, Gustav}!zzzSchnitzler, Olga@\emph{von Olga Schnitzler}!1901-06-171@{17. 6. 1901}|)be}\briefempfaengerindex{Schwarzkopf, Gustav@\textsc{Schwarzkopf, Gustav}!zzzSchnitzler, Arthur@\emph{von Arthur Schnitzler}!1901-06-171@{17. 6. 1901}|)be}\mylabel{L04502h}
\begin{anhang}
\end{anhang}\newcommand{\dateiname}{L04502}\newcommand{\titel}{Arthur Schnitzler und Olga Gussmann an Gustav Schwarzkopf, 17. 6. 1901}\newcommand{\editorInnen}{Herausgegeben von Jahnke, SelmaMüller, Martin Anton}\input{../tex-inputs/latex-leseansicht-abspann}
